\documentclass[11pt,oneside]{book}

\usepackage[T1]{fontenc}
\usepackage[utf8]{inputenc}
\usepackage{lmodern}
\usepackage[margin=1in]{geometry}
\usepackage{amsmath,amssymb,amsthm,mathtools}
\usepackage{microtype}
\usepackage{booktabs,longtable,tabularx,array}
\usepackage{enumitem}
\usepackage[most]{tcolorbox}
\usepackage{xcolor}
\usepackage{hyperref}
\usepackage{fancyhdr}
\usepackage{titlesec}

\definecolor{Primary}{HTML}{1F3A5F}
\definecolor{Accent}{HTML}{2D6A8D}
\definecolor{SoftBlue}{HTML}{EEF5FB}
\definecolor{SoftGold}{HTML}{FFF8E1}
\definecolor{SoftGreen}{HTML}{EDF7ED}
\definecolor{SoftGray}{HTML}{F4F4F4}

\hypersetup{
  colorlinks=true,
  linkcolor=Primary,
  urlcolor=Accent,
  pdftitle={Ontario Intermediate Mathematics ABQ Study Textbook},
  pdfauthor={Campbe1xx/Math-Test-Prep Repository},
  pdfsubject={Ontario Grade 9/10 Mathematics ABQ Study Resource}
}

\pagestyle{fancy}
\fancyhf{}
\fancyhead[L]{\leftmark}
\fancyhead[R]{Ontario Intermediate Mathematics ABQ Study Textbook}
\fancyfoot[C]{\thepage}
\setlength{\headheight}{14pt}

\titleformat{\chapter}[display]
  {\normalfont\bfseries\color{Primary}}
  {\filleft\Huge\chaptername\ \thechapter}
  {1ex}
  {\filright\Huge}

\setlist[itemize]{topsep=0.4em,itemsep=0.3em}
\setlist[enumerate]{topsep=0.4em,itemsep=0.35em}

\newtcolorbox{summarybox}{
  breakable,
  colback=SoftBlue,
  colframe=Primary,
  title=Chapter Summary,
  fonttitle=\bfseries
}

\newtcolorbox{notebox}[1]{
  breakable,
  colback=SoftGold,
  colframe=Accent,
  title=#1,
  fonttitle=\bfseries
}

\newtcolorbox{exercisebox}[1]{
  breakable,
  colback=SoftGray,
  colframe=black!60,
  title=#1,
  fonttitle=\bfseries
}

\tcbset{
  definitionstyle/.style={breakable,colback=SoftBlue,colframe=Accent,fonttitle=\bfseries},
  examplestyle/.style={breakable,colback=white,colframe=Primary,fonttitle=\bfseries},
  theoremstyle/.style={breakable,colback=SoftGreen,colframe=green!45!black,fonttitle=\bfseries}
}

\newtcbtheorem[number within=chapter]{definition}{Definition}{definitionstyle}{def}
\newtcbtheorem[number within=chapter]{example}{Worked Example}{examplestyle}{ex}
\newtcbtheorem[number within=chapter]{theorembox}{Key Idea}{theoremstyle}{th}

\newcommand{\coursecode}[1]{\textbf{#1}}
\newcommand{\topicline}[1]{\medskip\noindent\textbf{Focus:} #1\par\medskip}
\newcommand{\answerline}{\par\smallskip\noindent\textit{Answer: }}

\begin{document}

\frontmatter

\begin{titlepage}
\centering
{\Huge\bfseries Ontario Intermediate Mathematics ABQ Study Textbook\par}
\vspace{1.5cm}
{\Large A textbook-style study package based on Ontario Grade 9 and Grade 10 mathematics\par}
\vspace{1cm}
{\large for teacher-candidates and learners preparing for the Intermediate Mathematics ABQ\par}
\vfill
\begin{tcolorbox}[width=0.9\textwidth,colback=SoftGold,colframe=Accent,title=Important Disclaimer,fonttitle=\bfseries]
This document is a \textbf{curriculum-aligned study package}. It is \textbf{not an official Trent University practice test}, and it does not claim Trent-specific weighting, timing, cut-scores, or question distribution. It reorganizes the repository's existing Ontario Intermediate Mathematics ABQ study content into a polished textbook-style format for review and practice.
\end{tcolorbox}
\vfill
\begin{tabular}{rl}
Source basis: & \texttt{Ontario-Intermediate-Math-ABQ-Study-Package.md} \\
Curriculum scope: & Ontario Grade 9 Mathematics \coursecode{MTH1W} strands B--F \\
& Ontario Grade 10 Academic Mathematics \coursecode{MPM2D} (2005), strands A--C
\end{tabular}
\vfill
{\large Campbe1xx/Math-Test-Prep Repository\par}
{\large \today\par}
\end{titlepage}

\chapter*{How to Use This Textbook}
\addcontentsline{toc}{chapter}{How to Use This Textbook}
This study text is designed for two overlapping audiences:
\begin{itemize}
  \item Grade 9 or Grade 10 Ontario mathematics learners who want a concise reference and practice resource.
  \item Teacher-candidates preparing for the Intermediate Mathematics ABQ who need a coherent review of the Ontario curriculum foundations named in the repository study package.
\end{itemize}
Use each chapter in the same pattern:
\begin{enumerate}
  \item Read the opening ideas and definitions.
  \item Study the worked examples carefully.
  \item Complete the chapter exercises without looking at the answers.
  \item Use the appendix answer key to check fluency, not just final answers.
  \item Revisit the curriculum overview if you need to connect a skill to Ontario expectations.
\end{enumerate}

\tableofcontents

\mainmatter

\chapter{Ontario Intermediate Mathematics ABQ Overview}
\section{Purpose and Scope}
The original repository package prepares readers for an entry test aligned to \textbf{Ontario Grades 9 and 10 mathematics}. This textbook retains that same organizing principle. The emphasis is not on memorizing isolated tricks, but on recovering the habits of mind expected in intermediate mathematics:
\begin{itemize}
  \item numerical fluency,
  \item algebraic reasoning,
  \item representation of relations,
  \item data-informed decision-making,
  \item geometric and spatial reasoning,
  \item financial literacy,
  \item quadratic thinking, and
  \item trigonometric modelling.
\end{itemize}

\section{Study Plans from the Existing Package}
\subsection*{Two-day compressed review}
\begin{itemize}
  \item Day 1: fluency repair, diagnostic review, Grade 9 number and algebra, and the foundations of quadratics.
  \item Day 2: data, financial mathematics, analytic geometry, trigonometry, and a cumulative mock review.
\end{itemize}

\subsection*{Three-day review}
\begin{itemize}
  \item Day 1: number sense, exponents, equations, and linear relations.
  \item Day 2: data, geometry, measurement, financial mathematics, and quadratic relations.
  \item Day 3: analytic geometry, trigonometry, mixed practice, and final readiness checking.
\end{itemize}

\section{Curriculum Alignment Overview}
\begin{longtable}{>{\raggedright\arraybackslash}p{0.17\textwidth} >{\raggedright\arraybackslash}p{0.22\textwidth} >{\raggedright\arraybackslash}p{0.5\textwidth}}
\toprule
Grade & Strand & Topics emphasized in this textbook \\
\midrule
\endhead
Grade 9 & Number & number sets, density, exponent laws, scientific notation, rational number operations, rates, ratios, percents \\
Grade 9 & Algebra & simplifying expressions, solving equations, coding representations, linear and non-linear relations, graph features \\
Grade 9 & Data & one-variable and two-variable data, modelling, scatter plots, bias, probability foundations \\
Grade 9 & Geometry and Measurement & measurement conversions, area, volume, right triangles, scale, spatial reasoning \\
Grade 9 & Financial Literacy & appreciation, depreciation, budgets, simple interest, quantitative comparison \\
Grade 10 Academic & Quadratic Relations & parabolas, transformations, factoring, completing the square, solving quadratic equations \\
Grade 10 Academic & Analytic Geometry & systems of linear equations, midpoint, distance, slope, circle equations, coordinate proofs \\
Grade 10 Academic & Trigonometry & similar triangles, right-triangle trigonometry, sine law, cosine law \\
\bottomrule
\end{longtable}

\chapter{Number Sense and Exponents}
\topicline{Number sets, rational numbers, exponent laws, scientific notation, ratio and percent reasoning}

\section{Number Systems and Density}
\begin{definition}{Number system}{numbersystems}
A \emph{number system} is a set of numbers with shared properties. Common sets include natural numbers, integers, rational numbers, and real numbers. Rational numbers are \emph{dense}: between any two rational numbers, there is always another rational number.
\end{definition}

\begin{theorembox}{Useful perspective}{density}
In Grade 9 mathematics, students move from seeing numbers as counting objects to seeing them as members of connected systems. This shift supports later work with algebra, slope, and real-valued measurement.
\end{theorembox}

\section{Exponent Laws}
For nonzero base $a$, the core exponent laws are:
\[
 a^m\cdot a^n=a^{m+n}, \qquad \frac{a^m}{a^n}=a^{m-n}, \qquad (a^m)^n=a^{mn}, \qquad a^0=1.
\]
These laws explain both integer exponents and scientific notation.

\begin{example}{Simplifying an expression}{exp-laws}
Simplify $\dfrac{2^5\cdot 2^{-3}}{2^2}$.
\tcblower
Combine exponents in the numerator first:
\[
2^5\cdot 2^{-3}=2^{5+(-3)}=2^2.
\]
Then divide:
\[
\frac{2^2}{2^2}=2^{2-2}=2^0=1.
\]
\answerline $1$
\end{example}

\section{Scientific Notation}
Scientific notation expresses a number in the form $a\times 10^n$, where $1\le a<10$ and $n$ is an integer. It is useful when comparing very large or very small quantities.

\begin{example}{Converting to scientific notation}{sci-notation}
Write $0.0062$ in scientific notation.
\tcblower
Move the decimal point three places to the right to obtain $6.2$. Because the original number is less than $1$, the exponent is negative:
\[
0.0062=6.2\times 10^{-3}.
\]
\answerline $6.2\times 10^{-3}$
\end{example}

\begin{exercisebox}{Exercises: Number sense and exponents}
\begin{enumerate}
  \item Classify each as natural, integer, rational, or irrational: $-3$, $\frac{7}{8}$, $\sqrt{2}$.
  \item Simplify $3^4\cdot 3^{-2}$.
  \item Simplify $\dfrac{x^5}{x^2}$ for $x\neq 0$.
  \item Write $54{,}000$ in scientific notation.
  \item A quantity grows from $80$ to $92$. Find the percent increase.
  \item Five pens cost \$11.25. What is the unit rate? What will eight pens cost at the same rate?
\end{enumerate}
\end{exercisebox}

\begin{summarybox}
Strong number sense shows up later in every chapter. If signs, fractions, and exponents feel shaky, pause here and rebuild that fluency before moving on.
\end{summarybox}

\chapter{Algebraic Thinking and Linear Relations}
\topicline{Expressions, equations, equivalent forms, graphs, slope, and multiple representations}

\section{Expressions and Equivalence}
Equivalent expressions name the same quantity in different ways. For example,
\[
-3(2x-5)+4x=-6x+15+4x=-2x+15.
\]
Equivalent forms matter because one form may be easier to solve, graph, or interpret than another.

\begin{definition}{Linear relation}{linearrelation}
A \emph{linear relation} has a constant first difference and a graph that forms a straight line. In the equation $y=mx+b$, the constant $m$ is the slope and $b$ is the $y$-intercept.
\end{definition}

\section{Solving Equations}
To solve a linear equation, isolate the variable while preserving equality by doing the same operation to both sides.

\begin{example}{Solving a multi-step equation}{lineq}
Solve $3(2x+1)=21$.
\tcblower
Distribute first:
\[
6x+3=21.
\]
Subtract $3$ from both sides:
\[
6x=18.
\]
Divide by $6$:
\[
x=3.
\]
\answerline $x=3$
\end{example}

\section{Slope and Intercepts}
For points $(x_1,y_1)$ and $(x_2,y_2)$ with $x_1\neq x_2$, the slope is
\[
 m=\frac{y_2-y_1}{x_2-x_1}.
\]
The sign of the slope tells whether the graph rises or falls from left to right.

\begin{example}{Finding the slope of a line}{slopeexample}
Find the slope through $(2,5)$ and $(7,-5)$.
\tcblower
\[
 m=\frac{-5-5}{7-2}=\frac{-10}{5}=-2.
\]
\answerline The slope is $-2$.
\end{example}

\section{Linear and Non-linear Patterns}
A table with constant first differences suggests a linear relation. Constant second differences suggest a quadratic relation, which becomes central in Grade 10.

\begin{exercisebox}{Exercises: Algebra and relations}
\begin{enumerate}
  \item Simplify $5x-2(3x-4)$.
  \item Solve $4x-7=17$.
  \item Solve $\frac{2}{3}x-5=9$.
  \item Find the slope of the line through $(1,4)$ and $(5,12)$.
  \item Determine the equation of the line with slope $-2$ and $y$-intercept $1$.
  \item Explain why the points $(1,2)$, $(3,6)$, and $(5,10)$ lie on a linear relation.
\end{enumerate}
\end{exercisebox}

\begin{summarybox}
Algebra is the bridge between arithmetic and modelling. By Grade 10, every quadratic, trigonometric, or financial context still depends on careful symbolic manipulation.
\end{summarybox}

\chapter{Data, Probability, and Mathematical Modelling}
\topicline{One-variable data, scatter plots, modelling decisions, theoretical and experimental probability}

\section{Reading and Representing Data}
Good data work requires more than calculation. You must decide what to measure, how to display it, and what limitations remain.

\begin{definition}{Theoretical and experimental probability}{probability}
\emph{Theoretical probability} is based on equally likely outcomes in a model. \emph{Experimental probability} is based on observed results from trials.
\end{definition}

\section{Measures of Centre}
Mean describes the arithmetic average, while median marks the middle value in an ordered list. Each gives a different view of the data.

\begin{example}{Comparing mean and median}{meanmedian}
For the data set $3,5,7,7,9$, find the mean and median.
\tcblower
The mean is
\[
\frac{3+5+7+7+9}{5}=\frac{31}{5}=6.2.
\]
The middle value in the ordered list is $7$, so the median is $7$.
\answerline Mean $=6.2$, median $=7$
\end{example}

\section{Scatter Plots and Modelling}
A scatter plot helps us describe association between two variables. A good mathematical model should fit the data reasonably well and respect the context that produced the data.

\section{Probability Foundations}
If a bag contains $5$ red cards and $7$ blue cards, then the probability of drawing a red card is
\[
\frac{5}{12}.
\]
If an experiment yields $18$ heads in $40$ flips, the experimental probability of heads is
\[
\frac{18}{40}=0.45.
\]

\begin{exercisebox}{Exercises: Data and probability}
\begin{enumerate}
  \item Identify one likely source of bias in an online poll about school attendance habits.
  \item Compute the mean and median of $4,4,6,10,11$.
  \item A fair six-sided die is rolled once. What is the probability of an odd result?
  \item A spinner lands on blue $9$ times out of $20$ trials. State the experimental probability of blue.
  \item Explain why a model with excellent numerical fit may still be weak in context.
\end{enumerate}
\end{exercisebox}

\begin{summarybox}
Ontario intermediate mathematics expects students to interpret data critically, not just compute with it. Always ask: what was measured, how was it measured, and what decision follows from the result?
\end{summarybox}

\chapter{Geometry, Measurement, and Coordinate Geometry}
\topicline{Measurement conversion, area and volume, Pythagorean relationships, midpoint, distance, slope, and geometric verification}

\section{Measurement and Spatial Reasoning}
Intermediate mathematics connects formulas to meaning. Units, diagrams, and structure matter as much as arithmetic.

\begin{theorembox}{Pythagorean theorem}{pythagorean}
In a right triangle with legs $a$ and $b$ and hypotenuse $c$,
\[
a^2+b^2=c^2.
\]
\end{theorembox}

\begin{example}{A right-triangle calculation}{righttriangle}
A right triangle has hypotenuse $13$ and one leg $5$. Find the other leg.
\tcblower
Use the Pythagorean theorem:
\[
5^2+b^2=13^2.
\]
So
\[
25+b^2=169, \qquad b^2=144, \qquad b=12.
\]
\answerline The missing leg is $12$.
\end{example}

\section{Coordinate Tools}
In Grade 10 analytic geometry, three formulas organize much of the work:
\[
\begin{aligned}
\text{slope } &m=\frac{y_2-y_1}{x_2-x_1},\\
\text{midpoint } &\left(\frac{x_1+x_2}{2},\frac{y_1+y_2}{2}\right),\\
\text{distance } &d=\sqrt{(x_2-x_1)^2+(y_2-y_1)^2}.
\end{aligned}
\]

\begin{example}{Midpoint and distance}{midpointdistance}
For $A(-4,6)$ and $B(8,-2)$, find the midpoint.
\tcblower
\[
\left(\frac{-4+8}{2},\frac{6+(-2)}{2}\right)=\left(\frac{4}{2},\frac{4}{2}\right)=(2,2).
\]
\answerline Midpoint $(2,2)$
\end{example}

\section{Geometric Verification in the Coordinate Plane}
Coordinate geometry allows us to justify that sides are parallel, perpendicular, equal in length, or bisected. This is a natural Grade 9 to Grade 10 transition: visual geometry becomes algebraic proof.

\begin{exercisebox}{Exercises: Geometry and coordinate geometry}
\begin{enumerate}
  \item Convert $2.4$ m to cm.
  \item Find the area of a triangle with base $10$ cm and height $7$ cm.
  \item Find the surface area formula you would use for a cylinder of radius $3$ and height $10$.
  \item Find the midpoint of $(2,1)$ and $(8,10)$.
  \item Decide whether slopes $3$ and $-\frac{1}{3}$ correspond to perpendicular lines.
  \item State one coordinate fact that would help prove a quadrilateral is a rectangle.
\end{enumerate}
\end{exercisebox}

\begin{summarybox}
Unit discipline and diagram discipline matter. Many errors that look algebraic are really geometric misunderstandings caused by poor sketches or ignored units.
\end{summarybox}

\chapter{Financial Mathematics}
\topicline{Percent change, appreciation, depreciation, simple interest, budgets, and comparison of options}

\section{Percent Change in Context}
Financial literacy in Ontario mathematics is not separate from algebra. Growth, decay, and comparison all rely on proportional reasoning.

\begin{definition}{Appreciation and depreciation}{appdep}
\emph{Appreciation} is an increase in value over time. \emph{Depreciation} is a decrease in value over time. In many Grade 9 contexts, a one-step change can be modelled as
\[
\text{new value}=\text{original value}\times(1\pm r),
\]
where $r$ is the rate written as a decimal.
\end{definition}

\begin{example}{A depreciation calculation}{depreciationex}
A laptop worth \$1200 depreciates by $12\%$ in one year. Find its value after one year.
\tcblower
A decrease of $12\%$ means multiplying by $0.88$:
\[
1200(0.88)=1056.
\]
\answerline The new value is \$1056.
\end{example}

\section{Simple Interest}
Simple interest is often modelled by
\[
I=Prt,
\]
where $P$ is principal, $r$ is the annual rate, and $t$ is the time in years.

\section{Budgets and Choice}
When income changes, a budget must be revised logically. Mathematics supports the comparison, but the decision still requires clear priorities.

\begin{exercisebox}{Exercises: Financial mathematics}
\begin{enumerate}
  \item Find the amount of a $15\%$ discount on \$240.
  \item A value rises from \$500 to \$560. Find the percent increase.
  \item Compute the simple interest on \$900 at $4\%$ for $3$ years.
  \item State two reasonable budget adjustments if monthly income drops by $8\%$.
\end{enumerate}
\end{exercisebox}

\begin{summarybox}
The strongest financial solutions explain both the calculation and the decision. A correct number with no interpretation is incomplete communication.
\end{summarybox}

\chapter{Quadratic Relations}
\topicline{Parabolas, transformations, factoring, completing the square, graph features, and solution methods}

\section{Recognizing a Quadratic Relation}
Quadratic relations often appear in standard form,
\[
y=ax^2+bx+c,
\]
and their tables show constant second differences. Their graphs are parabolas.

\begin{definition}{Vertex form}{vertexform}
A quadratic relation in \emph{vertex form} is written as
\[
y=a(x-h)^2+k,
\]
where $(h,k)$ is the vertex and the sign of $a$ determines whether the parabola opens upward or downward.
\end{definition}

\section{Factoring and Zeros}
Factoring reveals intercepts and supports efficient solving.

\begin{example}{Factoring and solving}{factorquad}
Solve $x^2-7x+12=0$.
\tcblower
Factor the trinomial:
\[
x^2-7x+12=(x-3)(x-4).
\]
Set each factor equal to zero:
\[
x-3=0 \quad \text{or} \quad x-4=0.
\]
So the solutions are $x=3$ and $x=4$.
\answerline $x=3$ or $x=4$
\end{example}

\section{Completing the Square}
Completing the square rewrites a quadratic in a form that reveals its vertex.

\begin{example}{From standard form to vertex form}{cts}
Write $y=x^2+6x+4$ in vertex form.
\tcblower
Group the first two terms and complete the square:
\[
y=(x^2+6x+9)-9+4=(x+3)^2-5.
\]
So the vertex form is
\[
y=(x+3)^2-5.
\]
\answerline $y=(x+3)^2-5$
\end{example}

\section{Graph Features}
From a graph or equation, identify:
\begin{itemize}
  \item the vertex,
  \item the axis of symmetry,
  \item the zeros or $x$-intercepts,
  \item the $y$-intercept, and
  \item the maximum or minimum value.
\end{itemize}

\begin{exercisebox}{Exercises: Quadratic relations}
\begin{enumerate}
  \item Expand $(x-5)^2$.
  \item Factor $x^2-9$.
  \item Solve $x^2-9=0$.
  \item State the axis of symmetry of $y=2x^2-8x+1$.
  \item State the vertex of $y=(x-3)^2-5$.
  \item Does $y=-2(x+1)^2+7$ open upward or downward?
\end{enumerate}
\end{exercisebox}

\begin{summarybox}
Quadratics are the major conceptual jump in the repository study package. Success depends on choosing the right representation: table, graph, factored form, standard form, or vertex form.
\end{summarybox}

\chapter{Analytic Geometry and Systems}
\topicline{Solving systems, comparing lines, coordinate arguments, and circle ideas}

\section{Systems of Linear Equations}
A system is solved by finding the ordered pair that satisfies both equations. Graphically, this is the point of intersection.

\begin{example}{Solving a system by elimination}{systemex}
Solve
\[
\begin{aligned}
x+y&=10,\\
x-y&=4.
\end{aligned}
\]
\tcblower
Add the equations:
\[
2x=14,
\]
so $x=7$. Substitute into $x+y=10$:
\[
7+y=10 \Rightarrow y=3.
\]
\answerline $(7,3)$
\end{example}

\section{Distance, Midpoint, and Circles}
The distance formula comes from the Pythagorean theorem. A circle centred at the origin with radius $r$ has equation
\[
x^2+y^2=r^2.
\]

\section{Coordinate Proof Ideas}
You can show two segments are parallel by equal slopes, perpendicular by negative reciprocal slopes, or congruent by equal distances.

\begin{exercisebox}{Exercises: Analytic geometry}
\begin{enumerate}
  \item Solve by substitution: $y=2x+1$ and $x+y=10$.
  \item Find the distance between $(2,1)$ and $(8,10)$.
  \item Write the equation of the circle centred at the origin with radius $5$.
  \item Name one coordinate condition that proves two lines are parallel.
\end{enumerate}
\end{exercisebox}

\begin{summarybox}
Analytic geometry rewards organization. Label points clearly, write formulas before substituting, and connect every calculation to a geometric claim.
\end{summarybox}

\chapter{Trigonometry and Similar Triangles}
\topicline{Similarity, right-triangle trigonometry, sine law, cosine law, and contextual problem solving}

\section{Similarity}
Similar triangles have equal corresponding angles and proportional corresponding sides. Similarity explains why trigonometric ratios stay constant for a fixed acute angle.

\begin{definition}{Primary trigonometric ratios}{trigratios}
In a right triangle with acute angle $\theta$,
\[
\sin\theta=\frac{\text{opposite}}{\text{hypotenuse}}, \qquad
\cos\theta=\frac{\text{adjacent}}{\text{hypotenuse}}, \qquad
\tan\theta=\frac{\text{opposite}}{\text{adjacent}}.
\]
\end{definition}

\begin{example}{Using tangent}{tanexample}
In a right triangle, the side opposite $\theta$ is $7$ and the adjacent side is $24$. Find $\theta$.
\tcblower
\[
\tan\theta=\frac{7}{24}.
\]
So
\[
\theta=\tan^{-1}\!\left(\frac{7}{24}\right)\approx 16.3^\circ.
\]
\answerline $\theta\approx 16.3^\circ$
\end{example}

\section{Sine Law and Cosine Law}
For an acute triangle with sides $a,b,c$ opposite angles $A,B,C$:
\[
\frac{a}{\sin A}=\frac{b}{\sin B}=\frac{c}{\sin C}
\]
and
\[
c^2=a^2+b^2-2ab\cos C.
\]
Use the sine law when you know an angle-side opposite pair. Use the cosine law when the included angle and two sides are known, or when all three sides are known.

\begin{exercisebox}{Exercises: Trigonometry}
\begin{enumerate}
  \item Explain the difference between similar and congruent triangles.
  \item In a right triangle, find the missing leg if the hypotenuse is $13$ and one leg is $5$.
  \item Use a trigonometric ratio to find an angle when the opposite side is $7$ and the adjacent side is $24$.
  \item State whether the sine law or cosine law is more natural for: sides $8$ and $10$ with included angle $60^\circ$.
  \item In the sine law setup $a=9$, $A=30^\circ$, $B=45^\circ$, write the equation used to find $b$.
\end{enumerate}
\end{exercisebox}

\begin{summarybox}
Trigonometry in the ABQ context is not only about calculator skill. It is about diagram choice, side-angle relationships, and selecting an efficient method.
\end{summarybox}

\chapter{Integrated Practice and Readiness}
\section{Mixed Review Set}
The original repository emphasizes mixed practice because real assessments rarely announce the topic before the question begins. Try the following without labels.

\begin{exercisebox}{Cumulative mixed practice}
\begin{enumerate}
  \item Simplify $-3(2x-5)+4x$.
  \item Convert $0.0375$ to a percent.
  \item Find the midpoint of $(-4,6)$ and $(8,-2)$.
  \item Solve $x^2-7x+12=0$.
  \item A die is rolled once. What is the probability of an odd number?
  \item A laptop worth \$1200 depreciates by $12\%$. Find the value after one year.
  \item For $y=(x-3)^2-5$, state the vertex.
  \item Solve the system $x+y=10$ and $x-y=4$.
  \item Find $\tan\theta$ if the opposite side is $7$ and the adjacent side is $24$.
  \item State one reason a scatter plot may show association without proving causation.
\end{enumerate}
\end{exercisebox}

\section{Final Readiness Checklist}
\begin{itemize}
  \item I can work fluently with signed numbers, fractions, rates, and percent.
  \item I can simplify expressions and solve linear equations accurately.
  \item I can identify slope, intercepts, midpoint, and distance.
  \item I can interpret one-variable and two-variable data.
  \item I can solve measurement and right-triangle problems with units.
  \item I can compare financial options using justified calculations.
  \item I can factor, solve, and interpret quadratic relations.
  \item I can solve systems and verify geometric properties in the coordinate plane.
  \item I can use trigonometric ratios, the sine law, and the cosine law appropriately.
\end{itemize}

\backmatter

\chapter{Answer Key}
\section*{Chapter 2: Number Sense and Exponents}
\begin{enumerate}
  \item $-3$ is an integer, $\frac{7}{8}$ is rational, $\sqrt{2}$ is irrational.
  \item $3^2=9$
  \item $x^3$
  \item $5.4\times 10^4$
  \item $15\%$
  \item Unit rate \$2.25 per pen; cost for eight pens is \$18.00.
\end{enumerate}

\section*{Chapter 3: Algebra and Relations}
\begin{enumerate}
  \item $-x+8$
  \item $x=6$
  \item $x=21$
  \item $2$
  \item $y=-2x+1$
  \item The slope between consecutive points is constant: $\frac{6-2}{3-1}=2$ and $\frac{10-6}{5-3}=2$.
\end{enumerate}

\section*{Chapter 4: Data and Probability}
\begin{enumerate}
  \item Sample answer: only people with internet access and motivation to respond may be represented.
  \item Mean $=7$, median $=6$.
  \item $\frac{3}{6}=\frac{1}{2}$
  \item $\frac{9}{20}$
  \item Sample answer: the model may ignore important context, bias, or unreasonable extrapolation.
\end{enumerate}

\section*{Chapter 5: Geometry and Coordinate Geometry}
\begin{enumerate}
  \item $240$ cm
  \item $35\text{ cm}^2$
  \item $2\pi r^2+2\pi rh$
  \item $(5,5.5)$
  \item Yes. They are negative reciprocals.
  \item Sample answer: adjacent sides are perpendicular and opposite sides are parallel.
\end{enumerate}

\section*{Chapter 6: Financial Mathematics}
\begin{enumerate}
  \item \$36
  \item $12\%$
  \item \$108
  \item Sample answers vary; examples include reducing discretionary spending and renegotiating recurring expenses.
\end{enumerate}

\section*{Chapter 7: Quadratic Relations}
\begin{enumerate}
  \item $x^2-10x+25$
  \item $(x-3)(x+3)$
  \item $x=\pm 3$
  \item $x=2$
  \item $(3,-5)$
  \item Downward
\end{enumerate}

\section*{Chapter 8: Analytic Geometry}
\begin{enumerate}
  \item $(3,7)$
  \item $\sqrt{117}=3\sqrt{13}$
  \item $x^2+y^2=25$
  \item Equal slopes
\end{enumerate}

\section*{Chapter 9: Trigonometry}
\begin{enumerate}
  \item Similar triangles have equal corresponding angles and proportional sides; congruent triangles are the same size and shape.
  \item $12$
  \item $\theta=\tan^{-1}(7/24)\approx 16.3^\circ$
  \item Cosine law
  \item $\dfrac{9}{\sin 30^\circ}=\dfrac{b}{\sin 45^\circ}$
\end{enumerate}

\section*{Chapter 10: Mixed Practice}
\begin{enumerate}
  \item $-2x+15$
  \item $3.75\%$
  \item $(2,2)$
  \item $x=3$ or $x=4$
  \item $\frac{1}{2}$
  \item \$1056
  \item $(3,-5)$
  \item $(7,3)$
  \item $\tan\theta=\frac{7}{24}$
  \item Sample answer: association alone does not establish a cause-and-effect relationship.
\end{enumerate}

\chapter{Appendix: Ontario Curriculum and Study Notes}
\section*{Curriculum reminder}
This textbook follows the same Ontario-aligned scope already documented in the repository study package:
\begin{itemize}
  \item Grade 9 Mathematics \coursecode{MTH1W}, strands B--F.
  \item Grade 10 Academic Mathematics \coursecode{MPM2D} (2005), strands A--C.
\end{itemize}

\section*{Why this remains a study package}
The repository content explicitly states that it is a study support resource rather than an official Trent University assessment. That distinction matters pedagogically and ethically: the goal is readiness through curriculum mastery, not prediction of test form.

\section*{Suggested compilation}
Compile with a standard PDF LaTeX workflow such as:
\begin{verbatim}
pdflatex Ontario-Intermediate-Math-ABQ-Textbook.tex
pdflatex Ontario-Intermediate-Math-ABQ-Textbook.tex
\end{verbatim}
A LaTeX build engine such as \texttt{latexmk -pdf} may also be used.

\end{document}
